% =========================================================
% EMENTAS DE LIBRAS I E LIBRAS II — RASCUNHO (15/10/2026)
% Retorno DIREN/CEPE, itens 17 e 18 + reunião de 15/10/2026 com a parecerista do CEPE.
%
% Regra de oferta (combinada em 15/10): Libras I e Libras II substituem imediatamente
% Espanhol — Ano 1 e Espanhol — Ano 2 quando houver ao menos 1 (um) estudante surdo
% matriculado na turma. Mesma CH, mesmo semestre, mesmo bloco (Formação Geral, Linguagens).
%
% Onde inserir no ementario_adm.tex:
%   - Libras I  -> logo após a TABELA EMENTA 5 (Espanhol — Ano 1)
%   - Libras II -> logo após a ementa Espanhol — Ano 2
% Base de conteúdo: UC "Libras" do PPC do Técnico em Tradução e Interpretação em
% Libras/Português Integrado (IFSC Câmpus Palhoça Bilíngue), adaptada para 80 h por ano.
%
% PENDENTE: (1) revisão por docente de Libras / Educação Especial;
%           (2) conferência dos títulos no acervo da Biblioteca do Câmpus Garopaba.
%
% ---------------------------------------------------------
% NOTA DE RODAPÉ PROPOSTA PARA AS MATRIZES DO 1º E DO 2º ANO
% (linhas "Espanhol -- Ano 1*" e "Espanhol -- Ano 2*"):
%
% * Língua Espanhola ou Língua Brasileira de Sinais (Libras), conforme art. 7º, §§ 1º e 2º,
%   das Diretrizes Curriculares dos Cursos Técnicos Integrados ao Ensino Médio do IFSC.
%   Havendo ao menos 1 (um) estudante surdo matriculado na turma, as unidades curriculares
%   Libras I (1º ano) e Libras II (2º ano) serão ofertadas em substituição imediata a
%   Espanhol — Ano 1 e Espanhol — Ano 2, com a mesma carga horária (80 h cada) e na mesma
%   organização semestral, sem alteração da carga horária total do curso.
% =========================================================

\clearpage
% ---------------------------------------------------------
% TABELA EMENTA 5A: Libras I — Ano 1 (oferta condicionada, ver nota *)
% ---------------------------------------------------------
\begin{xltabular}{\linewidth}{|X|p{2.5cm}|p{2.5cm}|}
\hline
\multirow{3}{=}{\textcolor{ifscgreen}{\textbf{Unidade Curricular:}}\\[3pt]\textbf{\large \revisao{Libras I — Ano 1*}}} & \multicolumn{2}{p{\dimexpr5cm+2\tabcolsep+\arrayrulewidth\relax}|}{\textcolor{ifscgreen}{\textbf{Semestre:}} \textbf{1º e 2º}} \\ \cline{2-3}
 & \textcolor{ifscgreen}{\textbf{CH EaD*:}} & \textcolor{ifscgreen}{\textbf{CH Total*:}} \\ \cline{2-3}
 & \textbf{00 h} & \textbf{80 h} \\ \hline
\endfirsthead
\multicolumn{3}{|p{\dimexpr\linewidth-2\tabcolsep-2\arrayrulewidth\relax}|}{\footnotesize \revisao{*Oferta em substituição a Espanhol — Ano 1 quando houver ao menos 1 (um) estudante surdo matriculado na turma.}} \\ \hline
\multicolumn{3}{|p{\dimexpr\linewidth-2\tabcolsep-2\arrayrulewidth\relax}|}{\textcolor{ifscgreen}{\textbf{Objetivos:}}} \\ \hline
\multicolumn{3}{|p{\dimexpr\linewidth-2\tabcolsep-2\arrayrulewidth\relax}|}{\revisao{
\begin{itemize}[leftmargin=*,noitemsep,topsep=2pt]
    \item Reconhecer a Língua Brasileira de Sinais como língua natural da comunidade surda brasileira, com estrutura gramatical própria e reconhecimento legal.
    \item Desenvolver competências comunicativas básicas em Libras por meio da compreensão e da produção de enunciados em situações cotidianas.
    \item Utilizar os parâmetros fonológicos da Libras (configuração de mão, ponto de articulação, movimento, orientação da palma e expressões não manuais) na produção de sinais.
    \item Empregar vocabulário contextualizado e estruturas frasais simples em diálogos do cotidiano escolar e social.
    \item Compreender aspectos da cultura e da identidade surdas, valorizando a diversidade linguística e a inclusão.
\end{itemize}
}} \\ \hline
\multicolumn{3}{|p{\dimexpr\linewidth-2\tabcolsep-2\arrayrulewidth\relax}|}{\textcolor{ifscgreen}{\textbf{Conteúdos:}}} \\ \hline
\multicolumn{3}{|p{\dimexpr\linewidth-2\tabcolsep-2\arrayrulewidth\relax}|}{\revisao{
Libras como língua: histórico da educação de surdos, mitos e reconhecimento legal (Lei nº 10.436/2002 e Decreto nº 5.626/2005). \newline
Parâmetros fonológicos da Libras e expressões não manuais. \newline
Alfabeto manual, números e soletração. \newline
Apresentação pessoal e formas de cumprimento; sinal pessoal. \newline
Pronomes pessoais, possessivos e interrogativos. \newline
Advérbios de tempo e introdução ao uso dos verbos em Libras. \newline
Vocabulário contextualizado: família, calendário, contexto escolar, cores, alimentos e bebidas, moradia, meios de transporte, esporte e lazer. \newline
Estruturas frasais simples: afirmativas, negativas e interrogativas. \newline
Cultura e identidade surdas: comunidade surda, artefatos culturais e acessibilidade. \newline
Abordagem de Temáticas Transversais: direitos humanos e inclusão da pessoa com deficiência (Lei nº 13.146/2015), respeito à diversidade linguística e prevenção de todas as formas de violência e discriminação.
}} \\ \hline
\multicolumn{3}{|p{\dimexpr\linewidth-2\tabcolsep-2\arrayrulewidth\relax}|}{\textcolor{ifscgreen}{\textbf{Estratégias de Ensino e Aprendizagem:}}} \\ \hline
\multicolumn{3}{|p{\dimexpr\linewidth-2\tabcolsep-2\arrayrulewidth\relax}|}{\textbf{Metodologia:} \revisao{Aulas expositivas e dialogadas conduzidas preferencialmente em Libras, com relação constante entre teoria e prática. Uso de recursos visuais (vídeos em Libras, imagens, dicionários digitais) e atividades práticas de conversação em duplas e em grupos, dramatizações e apresentações individuais. Interação com membros da comunidade surda sempre que possível. Os materiais e as atividades serão planejados de forma acessível para estudantes surdos e ouvintes. Sem carga horária EaD. Sem divisão de turma.} \newline\vspace{2pt}
\textbf{Avaliação:} \revisao{Avaliação processual e contínua, com base em atividades individuais e em grupo de compreensão e produção em Libras (diálogos, apresentações gravadas em vídeo e atividades de interpretação de enunciados simples). A recuperação de conteúdos ocorrerá ao longo da unidade curricular, por meio de novas atividades individuais.}} \\ \hline
\multicolumn{3}{|p{\dimexpr\linewidth-2\tabcolsep-2\arrayrulewidth\relax}|}{\textcolor{ifscgreen}{\textbf{Bibliografia Básica:}}} \\ \hline
\multicolumn{3}{|p{\dimexpr\linewidth-2\tabcolsep-2\arrayrulewidth\relax}|}{\revisao{QUADROS, Ronice Müller de. \textbf{Libras}. São Paulo: Parábola, 2019. \newline
GESSER, Audrei. \textbf{LIBRAS? que língua é essa?} crenças e preconceitos em torno da língua de sinais e da realidade surda. São Paulo: Parábola, 2009.}} \\ \hline
\multicolumn{3}{|p{\dimexpr\linewidth-2\tabcolsep-2\arrayrulewidth\relax}|}{\textcolor{ifscgreen}{\textbf{Bibliografia Complementar:}}} \\ \hline
\multicolumn{3}{|p{\dimexpr\linewidth-2\tabcolsep-2\arrayrulewidth\relax}|}{\revisao{CAPOVILLA, Fernando César; RAPHAEL, Walkiria Duarte; MAURICIO, Aline Cristina L. \textbf{Novo Deit-Libras}: dicionário enciclopédico ilustrado trilíngue da língua de sinais brasileira. 3. ed. São Paulo: Edusp, 2015. \newline
HONORA, Márcia; FRIZANCO, Mary Lopes Esteves. \textbf{Livro ilustrado de língua brasileira de sinais}: desvendando a comunicação usada pelas pessoas com surdez. São Paulo: Ciranda Cultural, 2009. \newline
STROBEL, Karin. \textbf{As imagens do outro sobre a cultura surda}. Florianópolis: Ed. da UFSC, 2008.}} \\ \hline
\end{xltabular}

\clearpage
% ---------------------------------------------------------
% TABELA EMENTA 22A: Libras II — Ano 2 (oferta condicionada, ver nota *)
% ---------------------------------------------------------
\begin{xltabular}{\linewidth}{|X|p{2.5cm}|p{2.5cm}|}
\hline
\multirow{3}{=}{\textcolor{ifscgreen}{\textbf{Unidade Curricular:}}\\[3pt]\textbf{\large \revisao{Libras II — Ano 2*}}} & \multicolumn{2}{p{\dimexpr5cm+2\tabcolsep+\arrayrulewidth\relax}|}{\textcolor{ifscgreen}{\textbf{Semestre:}} \textbf{3º e 4º}} \\ \cline{2-3}
 & \textcolor{ifscgreen}{\textbf{CH EaD*:}} & \textcolor{ifscgreen}{\textbf{CH Total*:}} \\ \cline{2-3}
 & \textbf{00 h} & \textbf{80 h} \\ \hline
\endfirsthead
\multicolumn{3}{|p{\dimexpr\linewidth-2\tabcolsep-2\arrayrulewidth\relax}|}{\footnotesize \revisao{*Oferta em substituição a Espanhol — Ano 2 quando houver ao menos 1 (um) estudante surdo matriculado na turma.}} \\ \hline
\multicolumn{3}{|p{\dimexpr\linewidth-2\tabcolsep-2\arrayrulewidth\relax}|}{\textcolor{ifscgreen}{\textbf{Objetivos:}}} \\ \hline
\multicolumn{3}{|p{\dimexpr\linewidth-2\tabcolsep-2\arrayrulewidth\relax}|}{\revisao{
\begin{itemize}[leftmargin=*,noitemsep,topsep=2pt]
    \item Ampliar as competências comunicativas em Libras para situações de interação social, escolar e profissional.
    \item Utilizar recursos gramaticais da Libras, como classificadores, verbos com concordância, uso do espaço de sinalização e marcação temporal.
    \item Compreender e produzir narrativas e textos sinalizados de diferentes gêneros, incluindo registros em vídeo.
    \item Aplicar a Libras em situações de atendimento ao público e de rotina organizacional, relacionando-a à atuação do técnico em Administração.
    \item Reconhecer os direitos linguísticos da pessoa surda e as práticas de acessibilidade comunicacional nas organizações.
\end{itemize}
}} \\ \hline
\multicolumn{3}{|p{\dimexpr\linewidth-2\tabcolsep-2\arrayrulewidth\relax}|}{\textcolor{ifscgreen}{\textbf{Conteúdos:}}} \\ \hline
\multicolumn{3}{|p{\dimexpr\linewidth-2\tabcolsep-2\arrayrulewidth\relax}|}{\revisao{
Revisão e ampliação do vocabulário de Libras I. \newline
Tipos de verbos em Libras: verbos simples, com concordância e espaciais. \newline
Uso do espaço de sinalização, referenciação e marcação temporal. \newline
Classificadores e descrição visual de pessoas, objetos e lugares. \newline
Expressões não manuais gramaticais e afetivas. \newline
Narrativas, diálogos e textos sinalizados; registro em vídeo e escrita de sinais (noções). \newline
Vocabulário do mundo do trabalho: profissões, atendimento ao público, rotinas administrativas, comércio, finanças, saúde e serviços públicos. \newline
Literatura surda e manifestações artísticas em Libras. \newline
Direitos linguísticos da pessoa surda e acessibilidade comunicacional nas organizações: o papel do tradutor e intérprete de Libras-Português. \newline
Abordagem de Temáticas Transversais: direitos humanos, inclusão no mundo do trabalho e combate ao capacitismo.
}} \\ \hline
\multicolumn{3}{|p{\dimexpr\linewidth-2\tabcolsep-2\arrayrulewidth\relax}|}{\textcolor{ifscgreen}{\textbf{Estratégias de Ensino e Aprendizagem:}}} \\ \hline
\multicolumn{3}{|p{\dimexpr\linewidth-2\tabcolsep-2\arrayrulewidth\relax}|}{\textbf{Metodologia:} \revisao{Aulas conduzidas preferencialmente em Libras, com atividades práticas de conversação, simulação de situações de atendimento ao público e de rotinas organizacionais, produção de narrativas e vídeos sinalizados, análise de produções da literatura surda e debates. Articulação com as unidades curriculares da Formação Técnica e com as Oficinas de Integração. Interação com membros da comunidade surda sempre que possível. Materiais e atividades acessíveis para estudantes surdos e ouvintes. Sem carga horária EaD. Sem divisão de turma.} \newline\vspace{2pt}
\textbf{Avaliação:} \revisao{Avaliação processual e contínua, com base em produções sinalizadas individuais e em grupo, simulações de atendimento e produção de vídeos em Libras. A recuperação de conteúdos ocorrerá ao longo da unidade curricular, por meio de novas atividades individuais.}} \\ \hline
\multicolumn{3}{|p{\dimexpr\linewidth-2\tabcolsep-2\arrayrulewidth\relax}|}{\textcolor{ifscgreen}{\textbf{Bibliografia Básica:}}} \\ \hline
\multicolumn{3}{|p{\dimexpr\linewidth-2\tabcolsep-2\arrayrulewidth\relax}|}{\revisao{QUADROS, Ronice Müller de; KARNOPP, Lodenir Becker. \textbf{Língua de sinais brasileira}: estudos linguísticos. Porto Alegre: Artmed, 2004. \newline
GESSER, Audrei. \textbf{O ouvinte e a surdez}: sobre ensinar e aprender a Libras. São Paulo: Parábola, 2012.}} \\ \hline
\multicolumn{3}{|p{\dimexpr\linewidth-2\tabcolsep-2\arrayrulewidth\relax}|}{\textcolor{ifscgreen}{\textbf{Bibliografia Complementar:}}} \\ \hline
\multicolumn{3}{|p{\dimexpr\linewidth-2\tabcolsep-2\arrayrulewidth\relax}|}{\revisao{CAPOVILLA, Fernando César; RAPHAEL, Walkiria Duarte; MAURICIO, Aline Cristina L. \textbf{Novo Deit-Libras}: dicionário enciclopédico ilustrado trilíngue da língua de sinais brasileira. 3. ed. São Paulo: Edusp, 2015. \newline
KARNOPP, Lodenir; KLEIN, Madalena; LUNARDI-LAZZARIN, Márcia (org.). \textbf{Cultura surda na contemporaneidade}: negociações, intercorrências e provocações. Canoas: Ed. ULBRA, 2011. \newline
STROBEL, Karin. \textbf{As imagens do outro sobre a cultura surda}. Florianópolis: Ed. da UFSC, 2008. \newline
BRASIL. \textbf{Decreto nº 5.626, de 22 de dezembro de 2005}. Regulamenta a Lei nº 10.436, de 24 de abril de 2002, que dispõe sobre a Língua Brasileira de Sinais -- Libras. Brasília, DF: Presidência da República, 2005.}} \\ \hline
\end{xltabular}
